% Part: first-order-logic
% Chapter: axiomatic-deduction
% Section: deduction-theorem-quantifiers

\documentclass[../../../include/open-logic-section]{subfiles}

\begin{document}

\olfileid[fr]{fol}{axd}{ddq}

\olsection{Le théorème de la déduction avec quantificateurs}

\begin{thm}[Théorème de la déduction]
\ollabel{thm:deduction-thm-q}
Si $\Gamma \cup \{!A\} \Proves !B$, alors
$\Gamma \Proves !A \lif !B$.
\end{thm}

\begin{proof}
Nous procédons de nouveau par récurrence sur la longueur de la
!!{derivation} de~$!B$ à partir de $\Gamma \cup \{!A\}$.

Le cas de base est identique à celui de la preuve du
\olref[ded]{thm:deduction-thm}.

Pour l'étape de récurrence, supposons de nouveau que la !!{derivation}
de~$!B$ à partir de $\Gamma \cup \{!A\}$ se termine par une étape~$!B$
justifiée par une règle d'inférence. Si cette règle est le modus
ponens, nous raisonnons comme dans la preuve du
\olref[ded]{thm:deduction-thm}. Si la règle est~\QR, supposons que
$!B \ident !C \lif \lforall[x][!D(x)]$ et qu'une !!{formula} de la
forme $!C \lif !D(a)$ figure plus tôt dans la !!{derivation}, où $a$
ne figure ni dans~$!C$, ni dans~$!A$, ni dans~$\Gamma$. Nous avons donc
\begin{align*}
  \Gamma \cup \{!A\} & \Proves !C \lif !D(a),\\
  \intertext{et l'hypothèse de récurrence s'applique; ainsi,}
    \Gamma & \Proves !A \lif (!C \lif !D(a)).\\
  \intertext{Or}
  & \Proves (!A \lif (!C \lif !D(a))) \lif ((!A \land !C) \lif !D(a));\\
  \intertext{le modus ponens donne donc}
  \Gamma & \Proves (!A \land !C) \lif !D(a).\\
  \intertext{Comme la condition de fraîcheur de l'eigenvariable reste satisfaite, nous pouvons ajouter une étape justifiée par~\QR{} et obtenir}
    \Gamma & \Proves (!A \land !C) \lif \lforall[x][!D(x)].\\
  \intertext{Nous avons aussi}
  & \Proves ((!A \land !C) \lif \lforall[x][!D(x)])
    \lif (!A \lif (!C \lif \lforall[x][!D(x)])),\\
  \intertext{et par modus ponens,}
    \Gamma & \Proves !A \lif (!C \lif \lforall[x][!D(x)]).
\end{align*}
Autrement dit, $\Gamma \Proves !A \lif !B$.

Nous laissons en exercice le cas où $!B$ est justifiée par~\QR, mais
est de la forme $\lexists[x][!D(x)] \lif !C$.
\end{proof}

\begin{prob}
Achever la preuve du
\olref[fol][axd][ddq]{thm:deduction-thm-q}.
\end{prob}

\begin{editorial}
Dans la source, l'avant-dernière formule affichée manque d'une
parenthèse fermante. Surtout, la phrase qui suit le calcul conclut
$\Gamma \Proves !B$, alors que le calcul établit
$\Gamma \Proves !A \lif !B$, qui est précisément la conclusion du
théorème. La parenthèse et la conclusion ont été rétablies.
\end{editorial}

\end{document}
